Design Organisational Approval

Privileges: A DOA holder can
Perform design activities within the scope of approval
Have compliance documents accepted by the Agency without further verification
Perform activities independently from the Agency

We are most interested in the first rule, and for appliance we have to read:
nitial Airworthiness Regulation, including Part 21
Relevant section of Regulation: Articles 1(h) & 8
Annex I to Regulation (EC) No 748/2012 - Part 21 - Certification of aircraft and related products, parts and appliances, and of design and production organisations
Relevant section of Part 21: Subpart J - Design Organisation Approval

Relevant link once I find it: https://eur-lex.europa.eu/legal-content/EN/TXT/?uri=CELEX%3A02012R0748-20230825
We need to cover: 
-safety management element
-design assurance element
-safety risk management
-Overview of the entire logistic system, including maintenance, product and parts
- Agency a handbook that describes, directly or by cross reference, the organisation, its relevant policies, processes and procedures (we are delft??)
For the change we just need the following:
After the issue of a design organisation approval, each change to the design management system that is significant to the demonstration of compliance or to the airworthiness, operational suitability and environmental protection of the product, part or appliance shall be approved by the Agency before being implemented. The design organisation shall submit to the Agency an application for approval demonstrating, on the basis of the proposed changes to the handbook, that it will continue to comply with this Annex.

21.A.251 The terms of approval shall identify the types of design work, the categories of products, parts and appliances for which the design organisation holds a design organisation approval, and the functions and duties that the organisation is approved to perform with regard to the airworthiness, operational suitability and environmental characteristics of the products. For design organisation approvals covering type-certification or European Technical Standard Order (ETSO) authorisation for auxiliary power units (APUs), the terms of approval shall contain in addition the list of products or APUs. Those terms shall be issued as part of a design organisation approval.


\begin{comment}
REAT THE DOWN BULLSHIT
\end{comment}
If we apply for this DOA we will get certified so that we will have the capability of internally managing and approving ourselves. This means we as an entity that can approve our own design modification up to a certain extent without the need of an EASA direct approval. We become our own supervisors essentially. 

The benefits that we have after receiving the DOA is that we can internally classify a change as "Minor" or "Major". (KEEP THIS IN MIND) 

The process of documentation and proof of compliance involves the following steps: classifying the action, designing, performing compliance demonstration, validation as well as a statement of compliance, before issuing a STATEMENT OF APPROVAL.

21.A.263
In terms of design changes, we are allowed to perform modifications/repairs that fall under the "MINOR" category ourselves. We just have to follow the steps mentioned above once we get the DOA ourselves.

For modifications/repairs that fall under "Major" class, we can perform all the steps EXCEPT THE LAST ONE, we cannot issue a statement of approval.


Q: Would EASA allow us to classify the wing change as a "Minor" design consideration when we are establishing the DOA?

While reading the documentation for Design Organisational Approval, we found out that if we are approved for this category, we can internally design, validate and certify "minor" changes brought to the design. It is also stipulated that we have the capability to classify ourselves the class of the modification. My question is, when setting up the DOA, can we set it up in such a manner with EASA that the change of the wings or tail would be classified under "minor" for us? Would EASA allow that, and is it possible?



##Some fees and bullshit for this part, all of the following are taken from the follwing link: https://eur-lex.europa.eu/eli/reg_impl/2025/2347/oj
The appear to be the latest fees after which EASA follws,

To be noted the following:
2.   The fee to be paid by the applicant for a given certification task shall consist of one of the following:

(a)

a flat fee as set out in Part I of the Annex;

(b)

a variable fee.

3.   The variable fee referred to in paragraph 2, point (b), shall be established by multiplying the actual number of working hours by the hourly rate set out in Part II of the Annex.

Explanation: We have to pay the base fee fro certification and then the additional working hours from them:

Base fee: 
Hour fee: 250 euro


Question for EASA employ

Hello!
1. After looking over the DOA, we figured out that we still need a certification for our aircraft. As such, we were wondering how long it would take EASA to certify our plane, for example, under a CS-23 classification and how long it would take them to give us the DOA certification. They are charging on an hourly-based process
2. We are tring to also figure out the costs of the certification process and EASA charges on an hourly-based process. We were wondering approximately how long it would take them to: 
    a) Give the initial certification of the type of UAS
    b) The time it would take for EASA to review the validation and verification data from us, as holders of the DOA, we would have done all the testing and validation of the Major changes that we bring each time we build new wings  



    @article{pylon_load_analysis,
author = {Wang, Ding and Niu, Yaopeng and Guo, Xiangying},
year = {2025},
month = {11},
pages = {012031},
title = {Dynamic Analysis of Aircraft Engine-Pylon Coupled System},
volume = {3145},
journal = {Journal of Physics: Conference Series},
doi = {10.1088/1742-6596/3145/1/012031}
}

@article{dot,
  author  = {Singh, Avishreshth and Charak, Akhil and Biligiri, Krishna Prapoorna and Pandurangan, Venkataraman},
  title   = {Glass and carbon fiber reinforced polymer composite wastes in pervious concrete: Material characterization and lifecycle assessment},
  journal = {Resources, Conservation and Recycling},
  volume  = {182},
  pages   = {106304},
  year    = {2022},
  doi     = {10.1016/j.resconrec.2022.106304}
}

@article{Keiner2024,
  author  = {Keiner, Dominik and M{\"u}hlbauer, Andreas and Lopez, Gabriel and Koiranen, Tuomas and Breyer, Christian},
  title   = {Techno-economic assessment of atmospheric {CO2}-based carbon fibre production enabling negative emissions},
  journal = {Mitigation and Adaptation Strategies for Global Change},
  volume  = {28},
  number  = {52},
  year    = {2024},
  doi     = {10.1007/s11027-023-10090-5}
}

@book{roskamLandingGearSystems,
  author    = {Roskam, Jan},
  title     = {Airplane Design Part IV: Layout Design of Landing Gear and Systems},
  publisher = {DARcorporation},
  address   = {Lawrence, KS},
  year      = {2010},
  pages     = {75--102},
  isbn      = {978-1-884885-53-2}
}

@inproceedings{structural_loads_handbook,
  title={Structural Loads Handbook},
  author={Pedro Albuquerque},
  year={2011},
  url={https://api.semanticscholar.org/CorpusID:204770453},
  note = {(accessed on 18.05.2026)}
}

@article{ZHANG20233867,
title = {Magnesium research and applications: Past, present and future},
journal = {Journal of Magnesium and Alloys},
volume = {11},
number = {11},
pages = {3867-3895},
year = {2023},
note = {Magnesium and Its Alloys for Better Future - JMA 10th Anniversary},
issn = {2213-9567},
doi = {10.1016/j.jma.2023.11.007},
url = {https://www.sciencedirect.com/science/article/pii/S2213956723002694},
author = {Jianyue Zhang, Jiashi Miao, Nagasivamuni Balasubramani, Dae Hyun Cho, Thomas Avey, Chia-Yu Chang, and Alan A. Luo},
keywords = {Magnesium alloys, Structural applications, Lightweighting, Biomedical applications, Energy applications},
abstract = {As the lightest structural metal and one of the most abundant metallic elements on earth, magnesium (Mg) has been used as an “industrial metal”}
}